% test-026.tex -- comandos \spcoord y \spmQ
%
% Compilar:  lualatex-dev test-026.tex
% Validar:   show-pdf-tags --xml test-026.pdf | rnv-wrapp
%            verapdf --flavour ua2 --format text test-026.pdf
%
% Cada caso es un parrafo numerado, para poder referirse a el al
% escuchar el PDF. Los comentarios "doc:" indican la lectura que
% documenta el manual de usuario; en el resto, la lectura se revisa
% escuchando. Las entradas invalidas (que detienen la compilacion)
% no van aqui.
\DocumentMetadata
  {
    lang=es-CL,
    pdfversion=2.0,
    pdfstandard=ua-2,
    tagging=on,
    tagging-setup={math/setup={mathml-SE}},
  }
\documentclass{article}
\usepackage[spanish]{babel}
\usepackage{unicode-math}
\usepackage{spintent}
\usepackage{hyperref}
\hypersetup
  {
    pdfdisplaydoctitle = true,
    pdftitle           = {test-026: coordenadas y numeros mixtos},
  }
\pagestyle{empty}
\begin{document}

\section*{A. El comando spcoord}

% doc: «menos 3 coma 4»
1. Numérica: $\spcoord{-3, 4}$.

% doc: «1 coma 5 punto y coma 2 coma 5»
2. Decimales con punto y coma: $\spcoord{1,5; 2,5}$.

% doc: «a subíndice 1,2 coma b»
3. Coma dentro de llaves: $\spcoord{a_{1,2}, b}$.

% doc: «x subíndice 1 coma y subíndice 1»
4. Subíndices: $\spcoord{x_{1},y_{1}}$.

% doc: «1 mitad coma 3 cuartos»
5. Fracciones: $\spcoord{\frac{1}{2}, \frac{3}{4}}$.

% doc: sin lectura
6. Con silent: $\spcoord[silent=true]{a_{1,2}, b}$.

\section*{B. El comando spmQ}

% doc: «2 y la fracción con numerador 5 y denominador 7»
7. Número mixto: $\spmQ{2}{5}{7}$.

% doc: «menos 2 y la fracción con numerador 5 y denominador 7»
8. Entero negativo: $\spmQ{-2}{5}{7}$.

% doc: «2 enteros la fracción con numerador 5 y denominador 7»
9. Con join-word: $\spmQ[join-word=enteros]{2}{5}{7}$.

% doc: «se abren paréntesis 2 y la fracción con numerador 5 y denominador 7 se cierran paréntesis»
10. Con paréntesis: $\spmQ[wrap-parens=true]{2}{5}{7}$.

% doc: sin lectura
11. Con silent: $\spmQ[silent=true]{2}{5}{7}$.

% doc: «1 y la fracción con numerador 1 y denominador 2 coma 3»
% limitación conocida: la forma larga de la fracción de MathCAT no tiene cierre (se presta a ambigüedad)
12. En una coordenada: $\spcoord{\spmQ{1}{1}{2}, 3}$.

% doc: «x es igual a 3 y la fracción con numerador 1 y denominador 4 más 1»
% limitación conocida: la forma larga de la fracción de MathCAT no tiene cierre (se presta a ambigüedad)
13. Dentro de una fórmula: $x = \spmQ{3}{1}{4} + 1$.

\end{document}
